\documentclass[12pt,letter]{article}

% Esto es para poder escribir acentos directamente:
\usepackage[latin1]{inputenc}
% Esto es para que el LaTeX sepa que el texto est� en espa�ol:
\usepackage[spanish]{babel}
%m�rgenes
\usepackage[top=2.5cm,bottom=2.5cm,right=2cm,left=2cm]{geometry}

% Paquetes de la AMS:
\usepackage{amsmath, amsthm, amsfonts}

%multicolumna
\usepackage{multirow}

%hipertexto
\usepackage{hyperref}
\hypersetup{
  colorlinks=true,
  urlcolor=cyan
}

\usepackage{fancyhdr} %tipo de p�gina, LINDA biggrin.gif
%grafica
\usepackage{graphicx}

\setlength{\parindent}{0pt}

% Teoremas
%--------------------------------------------------------------------------
\newtheorem{thm}{Teorema}[section]
\newtheorem{cor}[thm]{Corolario}
\newtheorem{lem}[thm]{Lema}
\newtheorem{prop}[thm]{Proposici�n}
\theoremstyle{definition}
\newtheorem{defn}[thm]{Definici�n}
\theoremstyle{remark}
\newtheorem{rem}[thm]{Observaci�n}

% Atajos.
% Se pueden definir comandos nuevos para acortar cosas que se usan
% frecuentemente. Como ejemplo, aqu� se definen la R y la Z dobles que
% suelen representar a los conjuntos de n�meros reales y enteros.
%--------------------------------------------------------------------------

\def\RR{\mathbb{R}}
\def\ZZ{\mathbb{Z}}

% De la misma forma se pueden definir comandos con argumentos. Por
% ejemplo, aqu� definimos un comando para escribir el valor absoluto
% de algo m�s f�cilmente.
%--------------------------------------------------------------------------
\newcommand{\abs}[1]{\left\vert#1\right\vert}

% Operadores.
% Los operadores nuevos deben definirse como tales para que aparezcan
% correctamente. Como ejemplo definimos en jacobiano:
%--------------------------------------------------------------------------
\DeclareMathOperator{\Jac}{Jac}

%--------------------------------------------------------------------------
\pagestyle{fancy} %configuraci\'on para la pagina linda
\renewcommand{\sectionmark}[1]{\markboth{}{\thesection\ \ #1}} %cambios de comentarios
\lhead{\includegraphics[scale=0.15]{logoUSMC.eps}} %parte de arriba, izq
\chead{} %parte de arriba, centro
\rhead{\includegraphics[scale=0.10]{logoDI.png}} %parte de arriba, derecha, le agrega la marca del capitulo
\lfoot{} %parte de abajo, izq
\cfoot{} %parte de abajo, centro
\rfoot{\thepage} %parte de abajo, derecha, va el numero de la p�gina

%--------------------------------------------------------------------------
\begin{document}
\setcounter{page}{1}
\vspace{4cm}
\begin{center}
{\LARGE Ficha de Anteproyecto ``Feria de Software 2013''}\\
\today
\end{center}

\vspace{-0,2cm}
\section*{}

\begin{table}[ht]
\centering
\resizebox*{14cm}{!}{
  \begin{tabular}{|l|l|l|p{3cm}|}\hline
    Nombre del Proyecto: & \multicolumn{3}{|l|}{V.I.Pe.R. (\emph{Virtual Interactive Pet Robot})} \\ \hline
    Nombre pre-empresa:  & \multicolumn{3}{|l|}{Phyrex} \\ \hline
    Descripci\'on corta: & \multicolumn{3}{|p{14cm}|}{Simulador de mascota que explota caracteristicas de smartphones Android e interactua con robot Lego, para difusion del Dpto. de Informatica de la UTFSM.} \\ \hline
    Jefe de proyecto:    & Rodrigo Fr\'ias   & \href{mailto:rodrigo.frias@alumnos.usm.cl}{\nolinkurl{rodrigo.frias@alumnos.usm.cl}} & +56 9 83988257 \\ \hline
    \multirow{3}{*}{Equipo compuesto por:}& Celeste Bertin    & \href{mailto:celeste.bertin@alumnos.inf.utfsm.cl}{\nolinkurl{celeste.bertin@alumnos.usm.cl}} & +56 9 68410901 \\ \cline{2-4}
    & Patricio Carrasco & \href{mailto:patricio.carrascod@alumnos.usm.cl}{\nolinkurl{patricio.carrascod@alumnos.usm.cl}} & +56 9 50626689 \\ \cline{2-4}
    & Rocio Fernandez   & \href{mailto:rocio.fernandezu@alumnos.usm.cl}{\nolinkurl{rocio.fernandezu@alumnos.usm.cl}} & +56 9 62426549 \\ \hline
    Categor\'ia:         & \multicolumn{3}{|l|}{Educaci\'on \& Entretenimiento} \\ \hline
    Campus:              & \multicolumn{3}{|l|}{Santiago San Joaqu\'in} \\ \hline
  \end{tabular}
}
\end{table}

\vspace{0,5cm}
\section*{Resumen}
%P1: Descripcion problema \\ 36 (buscar incluir postulacion e informacion)
Se busca cautivar a potenciales futuros alumnos que ingresen a Ingenier�a Civil Inform�tica exhibiendo un producto que una diferentes aspectos y rubros demostrando la amplitud real que posee la carrera m�s all� de las nociones preconcebidas socialmente.

%P2: Solucion \\ 36
Ante esto, se propone una posible soluci\'on implementando una aplicaci\'on m\'ovil para Smartphones con Sistema Operativo Android, consistente en una mascota virtual cuyas acciones realizadas tendr\'an una reacci\'on en una mascota rob\'otica Lego Mindstorm, y viceversa.

%P3: Propuesta de Valor\\


%P4: Innovacion de la solucion \\ 26
Explotaci�n de funciones de un Smartphone android no convencionales para una mascota virtual, as� como la interacci�n entre la aplicaci�n virtual y la entidad rob�tica.
\end{document}
